% Conference poster, A0 portrait, built with tikzposter.
% Blocks stack inside columns; change the column widths to rebalance.
\documentclass[25pt,a0paper,portrait]{tikzposter}

\usepackage{amsmath}
\usepackage{booktabs}

\title{Poster Title in One Line}
\author{First Author, Second Author, Third Author}
\institute{Department, University \quad\textbullet\quad \texttt{you@example.edu}}

% Colours. Change "accent" to restyle the whole poster.
\definecolor{accent}{HTML}{1B5E20}
\definecolorstyle{Oberleaf}{}{
  \colorlet{backgroundcolor}{white}
  \colorlet{framecolor}{accent}
  \colorlet{titlefgcolor}{white}
  \colorlet{titlebgcolor}{accent}
  \colorlet{blocktitlebgcolor}{accent}
  \colorlet{blocktitlefgcolor}{white}
  \colorlet{blockbodybgcolor}{accent!6}
  \colorlet{blockbodyfgcolor}{black}
}
\usetheme{Simple}
\usecolorstyle{Oberleaf}

% A framed box that marks where a figure goes: \placeholder{height}
\newcommand{\placeholder}[1]{%
  \begin{center}
    \fbox{\parbox[c][#1][c]{0.9\linewidth}{\centering Figure placeholder}}
  \end{center}}

\begin{document}

\maketitle

\begin{columns}
  \column{0.5}

  \block{Motivation}{
    A poster is read from two metres away, in about a minute. State the
    question in two or three sentences and let the figures carry the rest.
    \begin{itemize}
      \item What is the open problem?
      \item Why does it matter now?
      \item What did you do about it?
    \end{itemize}
  }

  \block{Method}{
    Describe the approach in a few steps, with the one equation that matters:
    \begin{equation*}
      i\hbar\,\frac{\partial \psi}{\partial t} = \hat{H}\,\psi
    \end{equation*}
    \placeholder{22cm}
    \begin{enumerate}
      \item First step
      \item Second step
      \item Third step
    \end{enumerate}
  }

  \block{Setup}{
    \placeholder{20cm}
    One caption-length sentence saying what the figure shows.
  }

  \column{0.5}

  \block{Results}{
    \placeholder{26cm}
    \begin{center}
      \begin{tabular}{lcc}
        \toprule
        Method & Accuracy & Runtime \\
        \midrule
        Baseline & 91.2\% & 14.0 s \\
        Ours     & 94.7\% & 9.5 s \\
        \bottomrule
      \end{tabular}
    \end{center}
    Say what the reader should see in the figure and the table.
  }

  \block{Conclusion}{
    \begin{itemize}
      \item The one result a visitor should walk away with.
      \item What it changes.
      \item What comes next.
    \end{itemize}
  }

  \block{References}{
    \begin{enumerate}
      \item A. Author, \emph{Journal Name} \textbf{12}, 345 (2024).
      \item B. Author, \emph{Conference Proceedings}, 67--89 (2023).
    \end{enumerate}
  }

  \block{Acknowledgements}{
    Funding bodies and collaborators.
  }
\end{columns}

\end{document}
